% !TeX program = pdflatex
\documentclass[11pt,a4paper]{amsart}

\newcommand{\notelanguage}{finnish}
% Shared preamble for course notes and exercise sets.

\usepackage[T1]{fontenc}
\usepackage{lmodern}
\providecommand{\notelanguage}{english}
\usepackage[finnish,english]{babel}
\AtBeginDocument{\expandafter\selectlanguage\expandafter{\notelanguage}}

\newcommand{\notesfinnishlanguage}{finnish}
\ifx\notelanguage\notesfinnishlanguage
\newcommand{\notestheoremname}{Lause}
\newcommand{\noteslemmaname}{Lemma}
\newcommand{\notespropositionname}{Propositio}
\newcommand{\notescorollaryname}{Seurauslause}
\newcommand{\notesdefinitionname}{Määritelmä}
\newcommand{\notesexamplename}{Esimerkki}
\newcommand{\notesproblemname}{Tehtävä}
\newcommand{\notesexercisename}{Tehtävä}
\newcommand{\notesremarkname}{Huomautus}
\newcommand{\notessolutionname}{Ratkaisu}
\else
\newcommand{\notestheoremname}{Theorem}
\newcommand{\noteslemmaname}{Lemma}
\newcommand{\notespropositionname}{Proposition}
\newcommand{\notescorollaryname}{Corollary}
\newcommand{\notesdefinitionname}{Definition}
\newcommand{\notesexamplename}{Example}
\newcommand{\notesproblemname}{Problem}
\newcommand{\notesexercisename}{Exercise}
\newcommand{\notesremarkname}{Remark}
\newcommand{\notessolutionname}{Solution}
\fi

\usepackage[a4paper,margin=30mm]{geometry}
\usepackage{microtype}
\usepackage{graphicx}
\usepackage{enumitem}

\newcommand{\exerciseheading}[2]{%
  \par\noindent
  {\large\bfseries #1\par}
  \vspace{0.1em}
  \noindent{\large #2\par}
  \vspace{1.5em}
}

\usepackage{amsmath}
\usepackage{amssymb}
\usepackage{amsthm}
\usepackage{mathtools}
\usepackage{aliascnt}

% Linear systems: the argument gives the number of alignment pairs.
\newenvironment{linearsystem}[1]
  {\left\{\begin{alignedat}{#1}}
  {\end{alignedat}\right.}

\usepackage[hidelinks,pdfusetitle]{hyperref}

\numberwithin{equation}{section}
\setlist{itemsep=0.25em,topsep=0.5em}

\theoremstyle{plain}
\newtheorem{theorem}{\notestheoremname}[section]
\newaliascnt{lemma}{theorem}
\newtheorem{lemma}[lemma]{\noteslemmaname}
\aliascntresetthe{lemma}
\newaliascnt{proposition}{theorem}
\newtheorem{proposition}[proposition]{\notespropositionname}
\aliascntresetthe{proposition}
\newaliascnt{corollary}{theorem}
\newtheorem{corollary}[corollary]{\notescorollaryname}
\aliascntresetthe{corollary}

\theoremstyle{definition}
\newaliascnt{definition}{theorem}
\newtheorem{definition}[definition]{\notesdefinitionname}
\aliascntresetthe{definition}
\newaliascnt{example}{theorem}
\newtheorem{example}[example]{\notesexamplename}
\aliascntresetthe{example}
\newaliascnt{problem}{theorem}
\newtheorem{problem}[problem]{\notesproblemname}
\aliascntresetthe{problem}
\newtheorem{exercise}{\notesexercisename}

\theoremstyle{remark}
\newaliascnt{remark}{theorem}
\newtheorem{remark}[remark]{\notesremarkname}
\aliascntresetthe{remark}

\newenvironment{solution}
{
\begin{proof}[\notessolutionname]}
  {
\end{proof}}

\usepackage[nameinlink,noabbrev]{cleveref}

\crefname{theorem}{theorem}{theorems}
\crefname{lemma}{lemma}{lemmas}
\crefname{proposition}{proposition}{propositions}
\crefname{corollary}{corollary}{corollaries}
\crefname{definition}{definition}{definitions}
\crefname{example}{example}{examples}
\crefname{problem}{problem}{problems}
\crefname{exercise}{exercise}{exercises}
\crefname{remark}{remark}{remarks}

\ifx\notelanguage\notesfinnishlanguage
\crefname{theorem}{lause}{lauseet}
\crefname{lemma}{lemma}{lemmat}
\crefname{proposition}{propositio}{propositiot}
\crefname{corollary}{seurauslause}{seurauslauseet}
\crefname{definition}{määritelmä}{määritelmät}
\crefname{example}{esimerkki}{esimerkit}
\crefname{problem}{tehtävä}{tehtävät}
\crefname{exercise}{tehtävä}{tehtävät}
\crefname{remark}{huomautus}{huomautukset}
\fi

\newcommand{\NN}{\mathbb{N}}
\newcommand{\NNone}{\mathbb{N}_1}
\newcommand{\ZZ}{\mathbb{Z}}
\newcommand{\QQ}{\mathbb{Q}}
\newcommand{\RR}{\mathbb{R}}
\newcommand{\CC}{\mathbb{C}}
\newcommand{\PP}{\mathbb{P}}

\DeclareMathOperator{\im}{im}
\DeclareMathOperator{\rank}{rank}
\DeclareMathOperator{\syt}{syt}
\DeclarePairedDelimiter{\abs}{\lvert}{\rvert}
\DeclarePairedDelimiter{\norm}{\lVert}{\rVert}
\newcommand{\set}[1]{\left\lbrace#1\right\rbrace}
\newcommand{\maxset}[1]{\max\set{#1}}
\newcommand{\minset}[1]{\min\set{#1}}
\newcommand{\powerset}[1]{\mathcal{P}\!\left(#1\right)}


\title{Analyysi II\\Harjoitus 1}
\author{}
\date{}

\begin{document}

\exerciseheading{Analyysi II}{Harjoitus 1}

\begin{exercise}
  Olkoon \(f\colon \RR \to \RR\),
  \[
    f(x) =
    \begin{cases}
      1, & \text{kun } x < 2,\\
      2, & \text{kun } x \geqslant 2.
    \end{cases}
  \]
  Osoita jatkuvuuden \((\varepsilon,\delta)\)-määritelmän
  perusteella, että funktio \(f\)
  \begin{enumerate}[label=(\alph*), nosep]
    \item on jatkuva pistessä \(x = -2\),
    \item ei ole jatkuva pisteessä \(x = 2\).
  \end{enumerate}
\end{exercise}

\begin{exercise}
  Olkoon \(f\colon \RR \to \RR\),
  \[
    f(x) =
    \begin{cases}
      x, & \text{kun } x < 0,\\
      x^2, & \text{kun } x \geqslant 0.
    \end{cases}
  \]
  Osoita määritelmän perusteella, että funktio \(f\) on jatkuva origossa.
\end{exercise}

\begin{exercise}
  Todista oppikirjan Lause 4.1.7.

  \medskip
  \noindent\textsc{Lause 4.1.7:} Olkoon \((x_n)\) suppeneva lukujono, jonka
  raja-arvo on \(x \in (a,b)\), ja olkoon \(f\colon (a,b) \to \RR\)
  funktio, joka on jatkuva pisteessä \(x\). Jos
  \(\{x_n : n \in \NNone\} \subset (a,b)\), niin lukujono \((f(x_n))\)
  suppenee kohti lukua \(f(x)\).
\end{exercise}

\begin{exercise}
  Olkoon \(f\colon \RR \to \RR\) jatkuva.
  Määritellään funktio \(g\colon \RR \to \RR\) asettamalla, että
  \(g(x) = \abs{f(x)}\). Osoita, että \(g\) on jatkuva.
\end{exercise}

\begin{exercise}
  Olkoon \(f\colon \RR \setminus \set{0} \to \RR\),
  \(f(x) = \frac{1}{x}\). Osoita määritelmän perusteella, että funktio
  \(f\) on jatkuva.
\end{exercise}

\begin{exercise}
  \leavevmode\par
  \begin{enumerate}[label=(\alph*), nosep]
    \item Olkoon \(f\colon \RR \to \RR\) funktio, jolle
      \[
        f(x+y) = f(x) + f(y)
      \]
      kaikilla \(x,y \in \RR\). Osoita: jos funktio \(f\) on jatkuva
      origossa, niin \(f\) on jatkuva kaikkialla.
    \item Anna esimerkki funktiosta \(f\colon \RR \to \RR\), joka toteuttaa
      (a)-kohdan ehdot mutta ei saa jokaisessa pisteessä arvoa \(0\).
  \end{enumerate}
\end{exercise}

\end{document}
